% Ad Atticum 5.19 (ad-atticum-05-19)
% Source: https://raw.githubusercontent.com/PerseusDL/canonical-latinLit/master/data/phi0474/phi057/phi0474.phi057.perseus-lat1.xml
% Fetched by scripts/fetch_latin.py

% Dateline (Perseus): Scr. in castris ad Cybistra x K. Oct. a. 703 (51).

\ciceroLetterOpener{CICERO ATTICO salutem}

\ciceroSection{1} obsignaram iam epistulam eam, quam puto te modo perlegisse scriptam mea manu in qua omnia continentur, cum subito Apellae tabellarius a. d. xi Kal. Octobris septimo quadragesimo die Roma celeriter (hui tam longe!) mihi tuas litteras reddidit. ex quibus non dubito quin tu Pompeium exspectaris dum Arimino rediret et iam in Epirum profectus sis, magisque vereor, ut scribis, ne in Epiro sollicitus sis non minus quam nos hic sumus. de Atiliano nomine scripsi ad Philotimum ne appellaret Messallam.

\ciceroSection{2} itineris nostri famam ad te pervenisse laetor magisque laetabor si reliqua cognoris. filiolam tuam tibi †iam Romae† iucundam esse gaudeo, eamque quam numquam vidi tamen et amo et amabilem esse certo scio. etiam atque etiam vale.

\ciceroSection{3} de Patrone et tuis condiscipulis quae de parietinis in Melita laboravi ea tibi grata esse gaudeo. quod scribis libente te repulsam tulisse eum qui cum sororis tuae fili patruo certarit, magni amoris signum. itaque me etiam admonuisti ut gauderem; nam mihi in mentem non venerat. non credo inquis. ut libet; sed plane gaudeo, quoniam τὸ interest τοῦ φθονεῖν .
